We presented the first investigation of brain-LLM alignment in the context of divergent creative thinking, using fMRI data from participants performing the Alternative Uses Test alongside a non-creative control task. Our results reveal that alignment with creativity-relevant brain networks scales positively with model size and creative task performance during prompt processing, but weakens once models generate responses, suggesting that the generative process introduces representational dynamics that diverge from human creative cognition. We further show that post-training objectives have selective and interpretable effects: creativity-optimized training increases alignment with high-creativity neural responses while decreasing alignment with low-creativity ones, whereas reasoning-chain training produces the opposite pattern. Together, these findings suggest that the correspondence between LLM representations and the neural geometry of human creative thought is nuanced and sensitive to both the stage of processing examined and the nature of the training objective. We hope these results encourage further investigation into the neural plausibility of creative cognition in language models.